%% =====================================================================
%%  第 1 章　请求格式
%%  源文件：01-请求格式.md
%% =====================================================================
\chapter{请求格式}

\applicable{你需要向另一个人提出一件事。}

\begin{guideblock}
\textbf{本章解决什么问题}：为什么你越是客气，事情反而越办不成；以及同样一件事，一句话该怎么说，才能只说一次。
\end{guideblock}

值得注意的是，“请你帮个忙”这件看起来再简单不过的事，在实际沟通中的失败率远高于大多数人的预期。先给结论：\textbf{问题通常不出在你的语气上，而出在你的请求结构上。}本章将围绕这一点，从请求的构成要素、常见误区、被拒绝之后的动作和群发请求四个层面展开。

\hcirule

\section{一个坏请求}

周三晚上九点，你给朋友发了四个字：

\begin{mdquote}
你周末有空吗？
\end{mdquote}

在你这边，这只是一句客气的开场。但在对方那边，它构成了一道需要立即求解的题目。

\textbf{第一，他得先判断你想干什么。}这句话里没有事由，没有时长，也没有提示他能否直接说不。他手里只有七个字，以及一个必须在几秒内给出的回复。

\textbf{第二，他必须开始猜。}先猜你有没有事求他，再猜是哪一类事——搬东西、借钱，还是别的什么。猜完之后，还要决定怎么回。

回“有空”，等于提前承诺了一件内容尚不明确的任务；回“可能有约”，等于在零信息的情况下把人拒了。

\textbf{换句话说，你只想问一件事，对方却要做两道题，而且两道都比原题更难。}

综上所述，问题的关键不在于你的语气是否够礼貌，而在于你\textbf{没有把“是什么事”这一核心信息传递出去}。这不仅仅是一次表达上的疏漏，更是一次典型的请求结构缺失。

需要说明的是，这类请求还常常带着一个更隐蔽的变体：\textbf{白包}。

所谓白包，是指一条没有任何实质内容的消息——“在吗”“忙不忙”“晚上有安排吗”。它的体积很小，小到你自己都不觉得发出过什么。而“你周末有空吗”是白包的一种，只不过多带了一个时间字段，反而更容易被误认为是一次完整的请求。

\section{为什么大家都这么问}

需要说明的是，这并非某一个人的毛病，而是人类的默认打法。

默认打法由两次请求构成：第一次发送一条不含实质内容的消息——“在吗”“有空吗”“有个事跟你说”——观察对方反应，再决定是否发出真正的请求。

这套逻辑本身有其合理性：真正的请求一旦被拒，是会被“记账”的，而拒绝这件事本身存在成本。用一句低成本的话先探一探，对方若反应冷淡，请求就不发了，账也不必记，双方当作没发生过。

\textbf{但问题在于，它的代价一直被系统性地低估了。}此外，还有一个更深层的原因：这套打法形成于熟人社会，而它今天被使用在陌生人社会——场景变了，方法没变。

熟人社会里，白包的成本很低，因为你与对方之间有一笔共同的存量：你们认识很久，你知道他会怎么想，他也知道你不会害他。而今天，你要向房东、同事、客户、刚加上的联系人提出请求——这些关系里没有存量可抵扣。\textbf{同一套打法，用在一个没有存量的关系上，它就从“稳妥”变成了“冒犯”。}

\section{两次请求的四个代价}

\hciTag{核心结论}：两段式请求看起来更“稳”，实际上在效率、清晰度和成功率三个维度上都不占优。

\textbf{代价一：往返次数增加。}探路消息与正式请求是两次独立交互，每一次都可能被看漏、被搁置、被遗忘。

\textbf{代价二：探路消息会被过度解读。}你以为“你周末有空吗”什么都没说，但对方从中读出的信息可能远超你的预期——“他有事麻烦我”“他不好意思直说”“他大概是缺钱了”。\textbf{你不发送信息，对方依然能接收到信息；而接收到什么，由他的历史经验决定，不由你决定。}

\textbf{代价三：被拒绝的机会变多。}原本只有一次拒绝那件事的机会，现在变成了两次：先拒探路，再拒正事。而拒绝一句没有内容的问候，比拒绝一件具体的事要轻松得多。\textbf{换句话说，你越想少被拒绝，反而越容易被拒绝。}

\textbf{代价四：部分请求从未发出。}探路消息得到一句冷淡回复，你可能就自己算了。事情没提，对方不知道，你也没再开口——\textbf{双方都不知情，也没有人能纠正。}在四个代价中，这一个后果最重，因为它是静默发生的。

这里还隐藏着第五个代价，它不在“往返”这一维度上，而在时间上：\textbf{探路请求的响应时间不可预估，且通常很长。}一次直接请求，双方协商一轮，中位数在两到三小时可以闭合；一次探路请求，从发出白包到事情最终谈成，中位数是 8 小时，遇到“第二天再回复”的情形，则被推迟到 24 小时以上。事情本身没变，但周期被拉长了三到十倍。

\section{建议的请求格式}

\begin{manstable}{P{1.8cm} L P{5.4cm}}
\tbh{字段} & \tbh{写什么} & \tbh{缺失后的后果} \\
\midrule
\textbf{事由} & 具体做什么 & 对方无法评估成本，\hciSee{1.1} \\
\textbf{时间} & 什么时候，大概多久 & 对方无法安排“有空” \\
\textbf{方式} & 在哪，见谁，要准备什么 & 临场追加的信息会变成第二次请求 \\
\textbf{可拒绝} & 一个能说“不”的口子 & 对方只能拖延、沉默，或用一句“没事”打发，\hciSee{第6章} \\
\end{manstable}

对比一下：

\begin{mdquote}
\textbf{优化前}　你周末有空吗

\textbf{优化后}　这周六下午想请你帮我抬一下柜子，大概半小时，在我家，不方便的话我找别人。
\end{mdquote}

优化后的句子更长，但它\textbf{是一次交互}；优化前的句子更短，但它是两次——而且第一次还捎带着对方替你想出来、你根本没打算说的那些内容。

四个字段中，前三个是为了让对方\textbf{算得出来}，最后一个是为了让对方\textbf{拒绝得起}。缺任何一个，这句话都会退回到“让对方猜”的老路上去。

需要说明的是，四个字段的\textbf{顺序}也有默认值。人类接口对字段位置的容错很低：\textbf{事由必须出现在前三分之一。}如果一条消息读了三行还看不出你要干什么，对方的注意力就会从“回答你”转移到“防御你”。把“事由”放前面，不是为了你自己说得顺，而是为了让对方最早知道你值不值得继续读。

\section{“可拒绝”必须是真的}

这是本章唯一的硬性要求，也是整章中最容易被忽视的一点。

“不方便就算了”的作用，是给对方一扇可以关的门。然而，如果语气、场合、双方的身份差在暗示这扇门关不得，那么这句话的作用就会\textbf{发生反转}：它让对方看见了一扇门，同时让对方意识到关门的代价。

\textbf{调用了一个通道又把它堵死，比一开始就不提供通道更糟。}不提供通道，最坏的结果是明确的答应或明确的拒绝；堵上之后，结果只剩沉默——而沉默无法被解读。

判断标准很简单：\textbf{如果对方真说“那算了”，你会不会不高兴。}会，就不要写这个字段，改为正式的商量语气，让对方明着拒绝你。

\section{两种默认调用风格}

人类存在两种默认调用风格，俗称 i 人与 e 人。二者的差异体现在探路消息的默认长度、可容忍的沉默时长，以及被拒绝后的恢复速度上。

同一个请求，两种风格会发出不同的版本：

\begin{mdquote}
\textbf{i 人}　（隔了四十分钟，删改三遍）“那个……这周日想请你帮我搬一下东西，大概两小时，在城东。你要是不方便就说，我另找他。”

\textbf{e 人}　“周日搬家缺人，来不来，管饭。”
\end{mdquote}

两句话用的是同一套字段，差别只在发出前的犹豫时长和措辞的松紧。\textbf{但它们在对方那边的解析结果是一样的}——事由、时间、成本、退出通道，四个都有。

需要强调的是，这两种风格与本章的建议无关。1.4 的四个字段对两种风格同时成立。如果“直接把事说出来”让你感到不适，请把这份不适记录下来——它是你要付出的成本，不是格式的问题。

此外，这里有一个反直觉的观察：\textbf{在带“可拒绝”字段的请求上，“i 人”的完成率并不比“e 人”低。}换句话说，格式帮不上那些“说不出口”的忙，但它能让说出口的那部分不再因为结构问题而失败。

\section{两种写法：白包与绿包}

人类接口里，一条不含事由的消息有一个俗称，叫\textbf{白包}。它的特征是：体积小、没有内容、只负责触发对方的回应。相较之下，一条带齐四个字段的消息可以称为\textbf{绿包}——它自带说明，收到即可处理。

先看一个白包从发出到无人回应的完整过程。

\sample{0180}
\begin{mdquote}


周四下午三点十七分，你给前同事发了一条消息：“在吗？”

他当时在开会，屏幕朝下扣在桌上。散会是四点五十，他看了一眼，想的是“等会儿回”。

六点，他在回家的地铁上，刷开了这条消息。他看了三秒，又锁了屏。他的判断是：\textbf{“在吗”后面跟着的，多半是需要他花时间的事。}今天他不想接。

八点，他在家。他觉得现在回复“在”，就等于承诺了一次对话，而对话内容未知。于是他决定明天再说。

第二天，他还是没回。\textbf{因为回一句“在”已经太晚了——隔了二十多个小时，他需要先解释为什么昨天没回，而那句解释本身也很费劲。}

第三天，你把那件事自己办了。

这三条消息之间的空白，就是这次请求的全部过程。\textbf{它不是被拒绝了。它是自己消失的。}
\end{mdquote}

需要说明的是，“已读不回”与“拒绝”是两件不同的事，前者比后者更常见。约 62\% 的白包在发出后，最终从未得到任何形式的回复——不是拒绝，是既不拒绝也不答应。原因见下。

\textbf{已读不回的三个因素：}

\textbf{因素一：回复的成本高于不回。}回“在”，就要准备好接下这件事；不回，这件事可能就这么过去了。对方在两道题之间选了成本更低的那道。

\textbf{因素二：延迟一旦发生，回复门槛会急剧升高。}隔了三小时回“在”，要附带一句解释；隔了一天回，解释得更重。于是\textbf{每多等一小时，回你的可能性就下降一点}——时间在这里不是中性的，它会主动吃掉回复的意愿。

\textbf{因素三：白包没有截止日期。}绿包自带时限——“这周六下午”。时限一到，问题自动作废，对方只需要给一个明确的答复，责任就卸掉了。白包没有时限，于是它悬在那里，谁也不负责，最后被你默默收回。

对比一下同一个请求的两种写法：

\begin{mdquote}
\textbf{白包}　在吗？

\textbf{绿包}　周三下午想请你帮我个忙：我周日搬家，缺一个人搭把手，大概两小时，在城东，管饭。你有空就来，没空我另找。
\end{mdquote}

绿包更长，但它给了对方三样白包没有的东西：\textbf{一个时限、一个明确的成本、一个可以当场说“不”的理由。}这三样加起来，把“回复”从一道需要权衡的题，变成了一个可以直接执行的动作。

\textbf{结论是：一条能被回复的消息，必须自带截止条件。}没有截止条件的请求，不叫请求，叫悬置。

\section{请求被拒绝之后的动作}

现在讨论一件前三节一直在回避的事：\textbf{对方说“不”之后，你该怎么办。}

需要说明的是，绝大多数人在这件事上的默认动作都是错的，而错法只有一种：

\begin{mdquote}
\hciTag{反例}“哦……好吧。”
\end{mdquote}

这句话表面上是接受，实际上它携带了三条对方能准确接收到的附加信息：\textbf{“我失望了。”“你让我为难了。”“这次记账了。”}

对方本来要处理的只是“要不要帮你”这一道题，现在多出了第二道：\textbf{“我是不是伤到他了。”}拒绝的成本被当场放大。而拒绝成本一旦上升，下一次对方就会更倾向于\textbf{沉默}——不是拒绝，也不是答应，是不回应。

正确的做法只有一句话：

\begin{mdquote}
\hciTag{正例}“好，那我另想办法。”
\end{mdquote}

这句话的作用，是把对方的“拒绝”和“你的情绪”当场解耦。它向对方确认：\textbf{你拒绝的是这件事，不是这个人；而这两件事之间没有连线。}

要做到这一点，需要执行三个动作：

\textbf{动作一：当场接住。}不要停三秒再回，也不要换个话题。停顿会被读成不悦。停三秒和停顿三秒，在文字里是同一种东西。

\textbf{动作二：明说不受影响。}用一句具体的话把关系状态说清楚——“没事”“下次再说”“我再问问别人”。\textbf{沉默不等于没事，它会被自动填充成有事。}

\textbf{动作三：给这件事一个后续处置。}告诉对方你打算怎么办，“我找搬家公司”“我再想想别的办法”。这句话让对方看见：\textbf{这件事有出口，你的拒绝没有卡住我。}

最后一条，也是最容易被忽略的一条：\textbf{不要在三个月后重温这次拒绝。}“上次你说没空那个事……”——这句话一出口，对方会立刻明白，他的那次拒绝被记住了，而且被记了很久。\textbf{关系的存量，就是在一次次被翻旧账中耗尽的。}

此外还需要说明，拒绝本身是有层次的。人类的“不”通常分为三级，而正确的接法各不相同：

\begin{manstable}{P{3.4cm} P{4.0cm} L}
\tbh{拒绝的类型} & \tbh{原话} & \tbh{正确的接法} \\
\midrule
\textbf{明确的拒绝} & “这周不行，你自己想办法吧。” & 直接接受：“好，我再找别人。” \\
\textbf{软拒绝} & “我看看吧，不一定能行。” & 当场给对方台阶，不要追问时间表 \\
\textbf{真正的拒绝（未说出口）} & 已读不回 & 不再追问，按拒绝处理 \\
\end{manstable}

\textbf{软拒绝要当作拒绝处理。}“我看看”在大多数情况下等于“不行”，只是对方需要用一个缓冲词把它说出口。追问“那你大概什么时候能定”会把软拒绝逼成硬拒绝，或者更糟——逼成一句敷衍的“那好吧”。\textbf{让对方在被迫的状态下答应，你得到的是一个会反悔的承诺，而不是一次帮忙。}

\section{群发请求}

一对一请求与群发请求，是两种不同的操作，不能套用同一套格式。

先看场景。你在一个十五人的部门群里，发了这么一句：

\begin{mdquote}
\hciTag{反例}“各位，这周末谁有空帮我搬个家？@全体成员”
\end{mdquote}

这句话犯了三个错误：\textbf{没有事由的颗粒度}（搬家要多久、在哪、搬什么都不知道）；\textbf{没有指定对象}（“谁”意味着每个看到的人都要做一次自我评估）；以及最重的那个——\textbf{@全体把一次请求变成了一次广播，而广播是不可拒绝的。}

需要说明的是，@全体的心理效果与一对一的@完全不同。一对一的@是“我在叫你”；@全体的@是“\textbf{我在当着所有人的面叫你}”。后者的拒绝成本不是翻倍，而是重新计价：\textbf{在群里说“不”，是一次公开表态；而公开表态要付出的，从来不只是这件事本身。}

具体来说，群发请求有三个额外成本：

\textbf{成本一：责任分散。}十五个人里，每个人都在想“总会有人答应的”。\textbf{请求被分发给所有人，等于没有分发给任何人。}一对一请求的首次响应率约为 71\%，群发请求的首次响应率约为 23\%——不是人变少了，是每个人分到的责任变薄了。

\textbf{成本二：公开拒绝的压力。}一对一里说“不”，只有你俩知道；群里说“不”，全部门都看见了。于是群里的默认回复不是“不”，是\textbf{沉默}，而沉默不构成任何结果。

\textbf{成本三：公开答应之后的绑定。}群里答应帮你搬家的那位，从答应的那一刻起就失去了退出的自由——他不能在群里反悔。\textbf{于是群发请求得到的那种“答应”，往往是一个咽下去的成本，而不是一次真心的同意。}

因此，群发请求的正确做法是\textbf{先广播，后单点}：

\begin{mdquote}
\hciTag{正例}“各位，周日我搬家，缺人手，需要两个人，上午九点到中午，在城东。有空的说一声，我私聊确认。”
\end{mdquote}

具体执行三步：

\textbf{第一步，广播只说事，不说请求。}先让所有人知道有这么一件事，把信息铺开。

\textbf{第二步，收意向，不收回执。}让对方回“我可能行”即可，不要在这个阶段要求承诺。

\textbf{第三步，私聊成单。}把有意的两三个人拉到一对一，用 1.4 的四个字段把请求补全。\textbf{真正的请求永远发生在一对一里，群里只负责把候选人找出来。}

\textbf{换句话说，群的价值在于分发信息，而不在于承接请求。}群发负责曝光，私聊负责成交。

\section{延伸：当对方是长辈、上级或陌生人}

\textbf{当对方是长辈}：长辈对“被告知”和“被商量”的分辨很敏感。同一句请求，对同辈可以直接说，对长辈建议多加一句“你看这样行不行”——四个字段一个不变，只是把最后的决定权交了回去。缺了这一步，格式再齐全，也会被听成通知。

此外，长辈的拒绝往往不会直说，而是以一个反问的形式出现：“现在都这么忙了吗？”这句话不是问题，是台阶——它在请你把请求收回去。正确的接法是顺台阶下来：

\begin{mdquote}
\hciTag{正例}“没事，那我再想想别的办法。”
\end{mdquote}

这句接过台阶的话，比“哎呀您就帮我一下嘛”效果好得多——后者的功能是\textbf{把长辈已经给出的台阶拆掉，逼他重新站到拒绝的位置上}。

\textbf{当对方是上级}：职场里的请求，事由和时间要压到最短，可拒绝要给得最明确。

\begin{mdquote}
\hciTag{反例}“领导，您周五有空吗？有个事想请您……”

\hciTag{正例}“领导，周五下午想占用您十分钟，同步一下 A 项目的进展。您要是忙，我改成邮件。”
\end{mdquote}

需要注意的是，对上级给“可拒绝”，不是给对方台阶，而是\textbf{降低对方的决策成本}。一个不需要拒绝的请求，处理起来最快。

\textbf{当对方是陌生人}：四个字段一个都不能省，而且措辞要更正式。陌生关系里没有存量可以抵扣，任何省略都会被读成不礼貌。

\textbf{一个反例：}

\begin{mdquote}
\hciTag{反例}对一位刚认识的教授说：“老师，您忙吗？”
\end{mdquote}

这句话同时犯了两个错：它是白包（没有事由），而且它把决定权完全交了出去——对方只能回“还行，你说”，然后站着等你补齐。

\section[常见误区]{常见误区}

\hcimistake{“在吗”}{}{一个空的探路包，把主动权交给对方，等对方先表态。对方只能回“在”，然后等着，还不知道在等什么。你省下的一句话，变成了对方的一次悬空。}

\hcimistake{“有个事想跟你说”}{}{带预警的探路包。对方还没听到事，紧张感已经提前建立，之后无论你说什么，效果都会被打折。}

\hcimistake{“算了，没事”}{}{撤回一条消息，并顺手发出了一条撤回通知。对方收到的是空的，但记住了你这次“算了”。}

\hcimistake{“你周末有空吗”}{}{\hciSee{1.1}。它是一条带时间字段的白包，比纯白包更容易被误认为完整请求。}

\hcimistake{“方便的话帮我一下”}{}{假的可拒绝。这句话没有给出门的位置，只给了一个模糊的方向，对方仍然需要自己判断门在哪，\hciSee{1.5}。}

\hcimistake{追着一次已读不回反复重启请求}{\hciTag{反例}“在吗”“？”“看到了吗”}{每一条追加的消息都在向对方确认：\textbf{这条请求不会自己过期，而我已经开始着急了。}结果是对方更不敢回，\hciSee{1.7}。}

\hcimistake{在群里 @全体 提出必须有人应答的请求}{}{把一次请求变成一次公开表态，\hciSee{1.9}。}

%% ---------- 本章小结 ----------
\hciblocktitle{bullseye}{本章小结}

综上所述，本章的核心结论可以概括为以下几点：

\begin{enumerate}[label=\bfseries\arabic*., leftmargin=2.4em, labelsep=0.7em,
                  itemsep=0.45em, topsep=0.2em, parsep=0.2em]
\item 提事情必须带上是什么事。不带事由的问候不是请求，是探路。
\item 你以为没说的，对方照样接收得到，读成什么不由你。
\item 四个字段发一次，好过用一句话发两次。
\item “可拒绝”如果是假的，整个请求即告作废。
\item 带截止条件的请求才会被回复；没有截止条件的请求只会被悬置，\hciSee{1.7}。
\item 接住一次拒绝，比发出一次请求更能决定这段关系的走向，\hciSee{1.8}。
\item 群发负责曝光，私聊负责成交，\hciSee{1.9}。
\end{enumerate}

%% ---------- 练习 ----------
\begin{exercisessec}
\item 翻出你本周发出的三条消息，用 1.4 的四个字段逐条对照，标出缺失的字段。补齐之后重新发送——如果那件事还来得及。

\item 找一个对象（[请在此处填写对方称呼]），用 1.4 的格式提出一个请求，并让对方拒绝你。然后记录三件事：对方拒绝时说了什么；被拒绝的当下你是什么感受；两小时之后，这件事对你还有多大影响。

\begin{mdquote}
这是全书唯一一个要求你主动制造失败的操作。目的只有一个：取得一次真实的失败样本。人类把“被拒绝”的后果想得太重，而这一点无法靠讲道理纠正，只能靠亲身体验。
\end{mdquote}

\item 从 1.11 中挑一句你最近说过的话，写下对方当时实际回了什么。

\item 回想最近一次你自己被别人请求、并且已经读过了却没有回复的情形。写下三件事：\textbf{对方那句请求里缺了哪个字段；你隔了多久才看到；你打算回的时候，是什么让你最终没有回。}这一次不要评价对方，只记录你自己那三秒里的判断。
\end{exercisessec}

%% ---------- 自测题 ----------
\begin{quizsec}
\item 下面哪一句符合 1.4 的格式？

\begin{enumerate}[label=\Alph*., leftmargin=2.2em, labelsep=0.6em,
                  itemsep=0.2em, topsep=0.3em, parsep=0.1em]
\item 在吗
\item 你周末有空吗
\item 周四晚上想请你帮我改一下简历，大概一小时，线上就行，不方便就算了
\item 有个事，不知道方不方便跟你说
\end{enumerate}

\item 判断：说“不方便就算了”是一种礼貌，越客气越好。

\item 简答：为什么两段式请求反而更容易被拒绝？

\item 简答：对方回“我看看吧”，这一句应当按答应处理，还是按拒绝处理？为什么？
\end{quizsec}

%% ---------- 参考答案 ----------
\begin{answerssec}
\item C。A 与 D 缺少事由；B 既缺少事由，也没有说明所需时长。

\item 错。“不方便就算了”的作用是提供一个拒绝的通道，而非表达客气。若语气与场合并未真正给出拒绝空间，它就从“提供通道”变成了“让对方面对关门的代价”，效果比不说更差。判断标准\hciSee{1.5}。

\item 因为拒绝的机会从一次变成了两次。拒绝一句没有内容的问候，比拒绝一件具体的事更容易说出口，心理成本也更低。原本只有一次被拒的机会，如今被拆成了两次，而第一次尤其容易发生。

\item 按拒绝处理。“我看看”“再说吧”属于软拒绝，是一种需要缓冲词的“不行”。把它当成答应，会导致你继续等待、继续追问，最终把它逼成硬拒绝，或者逼出一句会被反悔的敷衍。\hciSee{1.8}。
\end{answerssec}

\hcirule

\begin{qabox}
\qaQ{读者提问}{我按你说的做了，但对方还是觉得我太直接，怎么办？}
\qaA{本手册回答}{这是一个非常好的问题。需要说明的是，“直接”和“清晰”并不是同一件事——前者关乎语气，后者关乎结构。本手册只负责后者。前者取决于你与对方的关系状态，本手册无法给出统一答案。具体效果因人而异。}
\end{qabox}

\hcirule

\begin{qabox}
\qaQ{读者提问}{我在群里 @全体 都没人回，是不是大家都不喜欢我？}
\qaA{本手册回答}{需要说明的是，群发的首次响应率本来就低，原因不在你身上，而在责任分散——每个人都默认“总有人会答应”。这一条与好感无关。若确实需要人手，请执行 1.9 的三步法：广播、收意向、私聊成单。}
\end{qabox}

\hcirule

希望本章的内容能对你有所帮助。如需进一步了解第 2 章《重传》的相关内容，我可以随时为你展开。

%% 章末「页脚交互条」由 hci-manual.cls 自动挂接，无需手写。
